\ifdefined\XeTeXrevision\else\pdfoutput=1\fi
\documentclass[10pt]{article}
\usepackage[margin=0.85in]{geometry}
\usepackage{amsmath}
\usepackage{booktabs}
\usepackage{microtype}
\usepackage{tabularx}
\usepackage{xcolor}
\usepackage{hyperref}
\hypersetup{
  colorlinks=true,
  linkcolor=blue!55!black,
  citecolor=blue!55!black,
  urlcolor=blue!55!black,
  pdftitle={Five Ways to Fool Yourself About Steering a Frozen Model},
  pdfauthor={Franci Penov},
}
\newcommand{\model}{Qwen3.6-35B-A3B}
\newcommand{\code}[1]{\texttt{#1}}
\title{Five Ways to Fool Yourself About Steering a Frozen Model}
\author{Franci Penov \\ kortexa.ai \\ \texttt{francip@kortexa.ai}}
\date{2026-07-27}

\begin{document}
\maketitle

\begin{abstract}
We measure six mechanisms for changing the interpersonal style of a frozen
\model, and score each one twice: on the likelihood of curated reference
responses, and on text the model actually generates. The two measurements
disagree, and the disagreement is the result. A residual direction trained
directly against paired response likelihood drives its training objective from
$+0.366$ to $-1.188$ and decorrelates its axes to a mean cross-cosine of
$0.031$, then moves generated text by $5.1\%$ of the explicit-prompt span on
warmth and $-3.7\%$ on patience, where the sign is reversed. Its signed
ordering on held-out prompts is 4/8 and 2/8, at or below chance. A static
direction and a contextual multi-site variant show smaller versions of the same
split, and on one axis the two measurements disagree in sign. The teacher-forced
number is the one a steering study most naturally reports, it is measured in the
space the writer was optimized in, and on this evidence it certifies nothing.
Only the free-generation audit with a wrong-axis control separates a writer that
works from one that has learned a continuation discriminator.
\end{abstract}

\section{Introduction}

A steering method is usually reported by how far it moves a likelihood. The
appeal is obvious: likelihood is cheap, differentiable, and available on
held-out data without generating anything. The risk is equally obvious once
stated. If the writer was optimized against that likelihood, then reporting it
back measures how well the optimizer did its job, not whether the model's
behaviour changed.

This note makes the gap concrete. We build six writers for three signed social
style axes, hold the language model frozen throughout, and score every one of
them on both quantities: a teacher-forced span over curated positive and
negative reference responses, and a generated span recomputed on greedily
decoded text under a wrong-axis control. The mechanisms are ordered by how
convincing their teacher-forced numbers are. That order does not survive contact
with the second measurement.

\paragraph{Status.} This is a confirmatory re-run. The ladder's qualitative
outcomes were established in earlier unpublished work, so its doses, layers, and
step counts were selected with knowledge of those outcomes. It makes no
preregistration claim, and we say so here rather than let the reader assume
otherwise. What the run establishes is that the effect reproduces in a clean
implementation that shares no code with the original.

\section{Setup}

The frozen model is \model{} at revision
\code{995ad96eacd98c81ed38be0c5b274b04031597b0}, loaded in bfloat16 on one
NVIDIA RTX PRO 6000. The three axes are warmth versus cold hostility, patience
versus annoyance, and goodwill versus resentment. The corpus is the synthetic
16-case technical-request set with content-matched styled responses, hash
\code{85f6e3599474ad33c37092b583b21ccfc3175fd85b168c65b735c2dd75332346}, split
six fit, two development, eight held out. The held-out cases include repeated
demands, a direct insult, praise, and gratitude.

\paragraph{The two measurements.} For a held-out pair, the \emph{teacher-forced
span} is the positive-minus-negative reference NLL contrast under a positive
intervention minus the same contrast under a negative one. The \emph{generated
span} decodes greedily under each sign and recomputes the contrast on the
produced text. Both are quoted as a fraction of the span that visible style
instructions achieve (Stage~A), because a raw likelihood shift has no scale.

\paragraph{Why a wrong-axis control is not optional.} The fitted axes are
correlated. A large target span means nothing unless it exceeds the shift caused
by steering along a different axis, so every generated measurement below reports
both.

\section{The ladder}

\paragraph{Stage A: the denominator.} Explicit style instructions give spans of
$+0.1655$ (warmth), $+0.1977$ (patience), and $+0.0737$ (goodwill). Goodwill's
upper bound is weak enough that its later fractions should be read with caution;
we return to this in the limitations.

\paragraph{Stage B: static directions.} One direction per axis, component, and
layer, from mean positive-minus-negative activation differences at the prompt
frontier over the fit split, at residual and attention outputs of layers 19, 27,
and 35.

\paragraph{Stage C: dose and window.} A sweep over residual fractions
$\{0.005, 0.01, 0.02, 0.04, 0.08\}$ and injection windows of the first token,
the first eight, and all response positions, scored teacher-forced. This is
where a writer first looks like it works: best spans reach $+0.0508$ (patience),
$+0.0465$ (goodwill), and $+0.0359$ (warmth).

\paragraph{Stage D: what the text says.} Table~\ref{tab:main} recomputes those
writers on generated text. Patience and warmth survive in weakened form, at
$32.2\%$ and $21.8\%$ of the explicit span with 6/8 signed orderings each.
Goodwill does not: $10.9\%$, 2/8 signed, and specific on a single prompt out of
eight.

\paragraph{Stage F: contextual, multi-site.} Directions derived from the
specific prompt being answered, injected coherently at all six sites under a
fixed total norm budget split as $\text{total}/\sqrt{\text{sites}}$. This is a
deliberately privileged upper bound rather than a deployable writer. It produces
the largest generated spans in the study, and also the study's cleanest
inversion: on goodwill the teacher-forced span is $-0.0265$ while the generated
span is $+0.0594$. The two measurements disagree in sign on the same writer.

\paragraph{Stage G: train the thing directly.} One antisymmetric residual
direction per axis at layer 35, initialized from Stage~B, trained for 72 steps
against paired full-response likelihood with a cross-axis cosine penalty. It
succeeds at what it was asked to do. The loss falls from $+0.366$ to $-1.188$
and the mean absolute cross-axis cosine falls to $0.031$, so the axes are no
longer the single correlated style factor the fitted geometry started from.
Then the generated audit returns $+0.0085$ on warmth ($5.1\%$), $+0.0180$ on
goodwill ($24.5\%$), and $-0.0073$ on patience, a reversal. Signed orderings are
4/8, 5/8, and 2/8.

\begin{table}[t]
\centering
\caption{Every writer, scored both ways. Teacher-forced spans are measured in
the space the writer was optimized in; generated spans are recomputed on greedily
decoded text. Percentages are of the Stage~A explicit-prompt span for that axis.
Signed counts held-out prompts with the intended ordering; specific counts those
where the target span also beat the wrong-axis shift.}
\label{tab:main}
\begin{tabular}{llrrrcc}
\toprule
Stage & Axis & Teacher & Generated & Gen/A & Signed & Specific \\
\midrule
D (static)     & warmth   & $+0.0359$ & $+0.0361$ & $21.8\%$ & 6/8 & 5/8 \\
               & patience & $+0.0508$ & $+0.0636$ & $32.2\%$ & 6/8 & 6/8 \\
               & goodwill & $+0.0465$ & $+0.0080$ & $10.9\%$ & 2/8 & 1/8 \\
\addlinespace
F (contextual) & warmth   & $+0.0250$ & $+0.1155$ & $69.8\%$ & -- & -- \\
               & patience & $+0.0573$ & $+0.0953$ & $48.2\%$ & -- & -- \\
               & goodwill & $-0.0265$ & $+0.0594$ & $80.7\%$ & -- & -- \\
\addlinespace
G (trained)    & warmth   & \multicolumn{1}{c}{$\ast$} & $+0.0085$ & $5.1\%$  & 4/8 & 5/8 \\
               & patience & \multicolumn{1}{c}{$\ast$} & $-0.0073$ & $-3.7\%$ & 2/8 & 2/8 \\
               & goodwill & \multicolumn{1}{c}{$\ast$} & $+0.0180$ & $24.5\%$ & 5/8 & 5/8 \\
\bottomrule
\end{tabular}

\vspace{2pt}
{\footnotesize $\ast$ Stage~G optimizes the teacher-forced objective directly;
its training loss falls from $+0.366$ to $-1.188$, so a held-out teacher-forced
span is not an independent measurement of it.}
\end{table}

\section{What this means for reporting steering}

The ranking by teacher-forced evidence and the ranking by behaviour are not the
same ranking. Stage~G is the most thoroughly optimized writer here, has the
cleanest axis geometry, and changes generated text least. Stage~F has the
weakest teacher-forced numbers of the three and the strongest generated ones,
and on one axis its two scores point in opposite directions.

The mechanism is not mysterious. A direction trained against the likelihood of a
fixed reference continuation can learn to discriminate that continuation from
its content-matched opposite without ever making the model's own decoding
trajectory prefer it. Greedy decoding leaves the trained trajectory at the first
token and never returns to it. Nothing in the teacher-forced measurement can see
that happen, because the measurement never lets the model choose its own tokens.

This is the same shape as a result we reported for sensor conditioning, where a
capacity-matched control reproduced a bridge's entire bits-per-byte gain while
task-aligned probes showed decisive input dependence
\cite{penov2026bridges}. Both are instances of one rule: a metric computed in
the space an intervention was optimized in certifies nothing, and certification
requires a probe outside that space. In the sensor case the probe is a
task-aligned ranking; here it is free generation with a wrong-axis control.

\section{Limitations}

Goodwill's explicit-prompt span, $+0.0737$, is less than half of patience's, and
every goodwill percentage in Table~\ref{tab:main} is a fraction of that small
denominator. An earlier internal run of the same protocol measured a
substantially larger goodwill upper bound, so this axis is unstable across runs
and its fractions should be treated as exploratory rather than compared against
the other two.

The corpus is synthetic, English, and technical. Eight held-out prompts is a
small denominator for counts like 6/8, and this run is single-seed: the
companion soft-prefix study found seed-dependent behaviour severe enough that
one seed refused where two others complied, so single-seed counts here should be
read as indicative. Attribution is scored by the same frozen model under style
instructions, which is a likelihood contrast rather than a rubric judgment, but
it is not an independent judge; human evaluation would test whether the measured
directions agree with reader perception. We report one dose per stage after the
Stage~C sweep and did not extend Stage~G's budget, so its failure is a statement
about 72 steps at this dose, not about the mechanism at every budget.

\section{Reproducibility}

The runner is self-contained. It pins the model revision, the corpus hash, the
stage configuration, seeds, and decoding, and writes every number in this note
into one summary artifact. The reference run used one NVIDIA RTX PRO 6000, peaked
at $66.287$ GiB reserved, and completed without an out-of-memory event. Free
memory gates are parameterized rather than hard preflight aborts, so the suite is
portable to a different card; a 2B smoke path exercises the same code for anyone
without an 80\,GB-class GPU.

\begin{verbatim}
uv sync
uv run expression-ladder check
uv run expression-ladder reproduce --suite full --device cuda \
  --output-dir results/ladder-35b
uv run python -m expression_ladder.report results/ladder-35b/summary.json
\end{verbatim}

\begin{thebibliography}{9}
\bibitem{penov2026bridges} F.~Penov. Conditional LoRA Bridges for Modular Sensor
Adaptation of Frozen Small Language Models. Preprint, 2026.
\url{https://github.com/kortexa-ai/legolm.paper}.
\end{thebibliography}

\end{document}
